\begin{frame}
\frametitle{Simulátory QtMips a QtRVSim}

\begin{itemize}
\item dříve se používalo prostředí MipsIt, které bylo ale velmi omezené
\item QtMips se pro výuku APO využívá od letního semestru 2019
\begin{itemize}
\item QtMIPS bylo vytvořeno jako diplomová práce Karla Kočího vedená Pavlem Píšou: \textit{Graphical CPU Simulator with Cache Visualization}, kterou si můžete prohlédnout \url{https://dspace.cvut.cz/bitstream/handle/10467/76764/F3-DP-2018-Koci-Karel-diploma.pdf}
\item Opravy, rozšíření a částečné přepracování Pavel Píša
\end{itemize}
\item v roce 2022 byl simulátor přepracován pro architekturu RISC-V, hlavní změny jsou od Jakuba Dupáka z jeho bakalářské práce z 2021 \textit{Graphical RISC-V Architecture Simulator - Memory Model and Project Management} odkaz: \url{https://dspace.cvut.cz/bitstream/handle/10467/94446/F3-BP-2021-Dupak-Jakub-thesis.pdf}
\item Alternativy:
\begin{itemize}
 \item RARS: Risc-V Assembler and Runtime Simulator -- IDE založené na systému MARS \url{https://github.com/TheThirdOne/rars}
 \item EduMIPS64: simulátor napsaný v Javě \url{https://www.edumips.org/}
\end{itemize}
\end{itemize}

\end{frame}


\begin{frame}
\frametitle{QtRVSim - Download}
\begin{itemize}
\item Windows, Linux, Mac \url{https://github.com/cvut/qtrvsim/releases}
\item Standardní součástí GNU/Linux distribucí Debian, Ubuntu, OpenSUSE (Factory), Arch, NuxOS a mnoha dalších
\item Naše nejnovější verze pro Ubuntu \url{https://launchpad.net/~qtrvsimteam/+archive/ubuntu/ppa}
\item Pro Suse, Fedora, Ubuntu, Debian  \url{https://software.opensuse.org/download.html?project=home\%3Ajdupak\&package=qtrvsim}
\item Online version \url{https://comparch.edu.cvut.cz/qtrvsim/app/}
\item LinuxDays 2019 QtMips talk – záznam interaktivní přednášky \url{https://youtu.be/fhcdYtpFsyw}, \url{https://pretalx.linuxdays.cz/2019/talk/EAYAGG/}
\item QtRvSim prezentace pro RISC-V International Academic and Training SIG \url{https://youtu.be/J6AcPZZ_ISg?t=12}
\end{itemize}
\end{frame}
